\documentclass[a4paper, 12pt]{article}
\usepackage[T2A]{fontenc}
\usepackage[english,russian]{babel}
\usepackage{amsmath}
\usepackage{amsthm}
\usepackage{amsfonts}
\usepackage{amssymb}
\usepackage{pgfplots}

\usepackage{enumitem}
\usepackage{longtable}
% Абзацы и геометрия страницы
\usepackage{indentfirst}
%\usepackage{geometry}
\usepackage[
    top=2cm,
    bottom=2cm,
    left=1cm,
    right=1cm
]{geometry}
\begin{document}


\begin{center}
    МИНИСТЕРСТВО НАУКИ И ВЫСШЕГО ОБРАЗОВАНИЯ РОССИЙСКОЙ ФЕДЕРАЦИИ

    ГОСУДАРСТВЕННОЕ ОБРАЗОВАТЕЛЬНОЕ УЧРЕЖДЕНИЕ ВЫСШЕГО ПРОФЕССИОНАЛЬНОГО ОБРАЗОВАНИЯ

    РЯЗАНСКИЙ ГОСУДАРСТВЕННЫЙ РАДИОТЕХНИЧЕСКИЙ УНИВЕРСИТЕТ

    ИМ. В. Ф. УТКИНА

    КАФЕДРА ВЫЧИСЛИТЕЛЬНОЙ И ПРИКЛАДНОЙ МАТЕМАТИКИ
\end{center}

\vspace{5 cm}

\begin{center}
    ЛАБОРАТОРНАЯ РАБОТА № 7

    ПО ДИСЦИПЛИНЕ "ТЕОРЕТИЧЕСКАЯ ИНФОРМАТИКА"

    НА ТЕМУ "ПРЕДСТАВЛЕНИЕ СИМВОЛОВ В ПАМЯТИ ЭВМ"

    ВАРИАНТ 14

\end{center}

\vspace{7 cm}

\begin{flushright}
    Выполнил:

    Студент группы 6413

    Косяков Д.А.

    Проверил

    Доцент кафедры ВПМ

    Никичкин Борис Владимирович
\end{flushright}

\begin{center}
    РЯЗАНЬ 2026
\end{center}

\pagebreak

\textbf{ Цель работы:} изучить представление символов русского алфавита в памяти ЭВМ в кодировках 866, 1251 и Unicode.

\textbf{ Задание:}
\begin{enumerate}[nosep]
    \item Записать коды символов фамилии, записанной заглавными буквами, в кодировке 866.
    Перевести коды в шестнадцатеричную систему счисления.
    \item Записать коды символов отчества, записанного первой заглавной и строчными буквами,
    в кодировке 1251. Перевести коды в шестнадцатеричную систему счисления.
    \item Записать слово, представленное последовательностью кодов в шестнадцатеричной системе
    счисления в кодировке Unicode.
\end{enumerate}

\textbf{ Выполнение задания:}

\textbf{1. Буквы русского алфавита и их коды.} Коды определены по таблице: в кодировке 866 буквы А--Я имеют коды 128--159,
а--п --- 160--175, р--я --- 224--239, Ё --- 240, ё --- 241; в кодировке 1251 буквы А--Я имеют коды 192--223,
а--я --- 224--255, Ё --- 168, ё --- 184; в Unicode буквы А--Я имеют коды 1040--1071, а--я --- 1072--1103,
Ё --- 1025, ё --- 1105. В скобках указаны шестнадцатеричные значения кодов.

{\small
\begin{longtable}{|c|c|c|c||c|c|c|c|}
\hline
Буква & 866 & 1251 & Unicode & Буква & 866 & 1251 & Unicode \\ \hline
\endhead
А & 128 (80) & 192 (C0) & 1040 (0410) & а & 160 (A0) & 224 (E0) & 1072 (0430) \\ \hline
Б & 129 (81) & 193 (C1) & 1041 (0411) & б & 161 (A1) & 225 (E1) & 1073 (0431) \\ \hline
В & 130 (82) & 194 (C2) & 1042 (0412) & в & 162 (A2) & 226 (E2) & 1074 (0432) \\ \hline
Г & 131 (83) & 195 (C3) & 1043 (0413) & г & 163 (A3) & 227 (E3) & 1075 (0433) \\ \hline
Д & 132 (84) & 196 (C4) & 1044 (0414) & д & 164 (A4) & 228 (E4) & 1076 (0434) \\ \hline
Е & 133 (85) & 197 (C5) & 1045 (0415) & е & 165 (A5) & 229 (E5) & 1077 (0435) \\ \hline
Ё & 240 (F0) & 168 (A8) & 1025 (0401) & ё & 241 (F1) & 184 (B8) & 1105 (0451) \\ \hline
Ж & 134 (86) & 198 (C6) & 1046 (0416) & ж & 166 (A6) & 230 (E6) & 1078 (0436) \\ \hline
З & 135 (87) & 199 (C7) & 1047 (0417) & з & 167 (A7) & 231 (E7) & 1079 (0437) \\ \hline
И & 136 (88) & 200 (C8) & 1048 (0418) & и & 168 (A8) & 232 (E8) & 1080 (0438) \\ \hline
Й & 137 (89) & 201 (C9) & 1049 (0419) & й & 169 (A9) & 233 (E9) & 1081 (0439) \\ \hline
К & 138 (8A) & 202 (CA) & 1050 (041A) & к & 170 (AA) & 234 (EA) & 1082 (043A) \\ \hline
Л & 139 (8B) & 203 (CB) & 1051 (041B) & л & 171 (AB) & 235 (EB) & 1083 (043B) \\ \hline
М & 140 (8C) & 204 (CC) & 1052 (041C) & м & 172 (AC) & 236 (EC) & 1084 (043C) \\ \hline
Н & 141 (8D) & 205 (CD) & 1053 (041D) & н & 173 (AD) & 237 (ED) & 1085 (043D) \\ \hline
О & 142 (8E) & 206 (CE) & 1054 (041E) & о & 174 (AE) & 238 (EE) & 1086 (043E) \\ \hline
П & 143 (8F) & 207 (CF) & 1055 (041F) & п & 175 (AF) & 239 (EF) & 1087 (043F) \\ \hline
Р & 144 (90) & 208 (D0) & 1056 (0420) & р & 224 (E0) & 240 (F0) & 1088 (0440) \\ \hline
С & 145 (91) & 209 (D1) & 1057 (0421) & с & 225 (E1) & 241 (F1) & 1089 (0441) \\ \hline
Т & 146 (92) & 210 (D2) & 1058 (0422) & т & 226 (E2) & 242 (F2) & 1090 (0442) \\ \hline
У & 147 (93) & 211 (D3) & 1059 (0423) & у & 227 (E3) & 243 (F3) & 1091 (0443) \\ \hline
Ф & 148 (94) & 212 (D4) & 1060 (0424) & ф & 228 (E4) & 244 (F4) & 1092 (0444) \\ \hline
Х & 149 (95) & 213 (D5) & 1061 (0425) & х & 229 (E5) & 245 (F5) & 1093 (0445) \\ \hline
Ц & 150 (96) & 214 (D6) & 1062 (0426) & ц & 230 (E6) & 246 (F6) & 1094 (0446) \\ \hline
Ч & 151 (97) & 215 (D7) & 1063 (0427) & ч & 231 (E7) & 247 (F7) & 1095 (0447) \\ \hline
Ш & 152 (98) & 216 (D8) & 1064 (0428) & ш & 232 (E8) & 248 (F8) & 1096 (0448) \\ \hline
Щ & 153 (99) & 217 (D9) & 1065 (0429) & щ & 233 (E9) & 249 (F9) & 1097 (0449) \\ \hline
Ъ & 154 (9A) & 218 (DA) & 1066 (042A) & ъ & 234 (EA) & 250 (FA) & 1098 (044A) \\ \hline
Ы & 155 (9B) & 219 (DB) & 1067 (042B) & ы & 235 (EB) & 251 (FB) & 1099 (044B) \\ \hline
Ь & 156 (9C) & 220 (DC) & 1068 (042C) & ь & 236 (EC) & 252 (FC) & 1100 (044C) \\ \hline
Э & 157 (9D) & 221 (DD) & 1069 (042D) & э & 237 (ED) & 253 (FD) & 1101 (044D) \\ \hline
Ю & 158 (9E) & 222 (DE) & 1070 (042E) & ю & 238 (EE) & 254 (FE) & 1102 (044E) \\ \hline
Я & 159 (9F) & 223 (DF) & 1071 (042F) & я & 239 (EF) & 255 (FF) & 1103 (044F) \\ \hline
\end{longtable}
}

\textbf{2. Фамилия, записанная заглавными буквами:} КОСЯКОВ.

\textbf{3. Буквы фамилии в кодировке 866.} Код заглавной буквы равен $128 + $ номер буквы в алфавите без Ё
(считая А нулевой). Перевод в шестнадцатеричную систему --- делением на 16
(частное --- старшая цифра, остаток --- младшая):

\begin{center}
\begin{tabular}{|c|c|c|c|}
\hline
Буква & Код (10) & Перевод & Код (16) \\ \hline
К & 138 & $138 : 16 = 8$, остаток $10 = A$ & 8A \\ \hline
О & 142 & $142 : 16 = 8$, остаток $14 = E$ & 8E \\ \hline
С & 145 & $145 : 16 = 9$, остаток $1$ & 91 \\ \hline
Я & 159 & $159 : 16 = 9$, остаток $15 = F$ & 9F \\ \hline
К & 138 & $138 : 16 = 8$, остаток $10 = A$ & 8A \\ \hline
О & 142 & $142 : 16 = 8$, остаток $14 = E$ & 8E \\ \hline
В & 130 & $130 : 16 = 8$, остаток $2$ & 82 \\ \hline
\end{tabular}
\end{center}

КОСЯКОВ в кодировке 866: 138 142 145 159 138 142 130 (10) = 8A 8E 91 9F 8A 8E 82 (16).

\textbf{4. Отчество, записанное с заглавной буквы:} Андреевич.

\textbf{5. Буквы отчества в кодировке 1251.} Код заглавной буквы равен $192 + $ номер буквы,
код строчной --- $224 + $ номер буквы (буква Ё не встречается):

\begin{center}
\begin{tabular}{|c|c|c|c|}
\hline
Буква & Код (10) & Перевод & Код (16) \\ \hline
А & 192 & $192 : 16 = 12 = C$, остаток $0$ & C0 \\ \hline
н & 237 & $237 : 16 = 14 = E$, остаток $13 = D$ & ED \\ \hline
д & 228 & $228 : 16 = 14 = E$, остаток $4$ & E4 \\ \hline
р & 240 & $240 : 16 = 15 = F$, остаток $0$ & F0 \\ \hline
е & 229 & $229 : 16 = 14 = E$, остаток $5$ & E5 \\ \hline
е & 229 & $229 : 16 = 14 = E$, остаток $5$ & E5 \\ \hline
в & 226 & $226 : 16 = 14 = E$, остаток $2$ & E2 \\ \hline
и & 232 & $232 : 16 = 14 = E$, остаток $8$ & E8 \\ \hline
ч & 247 & $247 : 16 = 15 = F$, остаток $7$ & F7 \\ \hline
\end{tabular}
\end{center}

Андреевич в кодировке 1251: 192 237 228 240 229 229 226 232 247 (10) = C0 ED E4 F0 E5 E5 E2 E8 F7 (16).

\textbf{6. Перевод шестнадцатеричных чисел в слово.} Дана последовательность кодов (вариант 14):
$$\mathrm{04370430043A043E043D0430}$$

В Unicode каждый символ русского алфавита занимает 2 байта, то есть 4 шестнадцатеричные цифры.
Разбиваем последовательность на группы по 4 цифры и переводим каждую в десятичную систему:

\begin{center}
\begin{tabular}{|c|c|c|}
\hline
Код (16) & Перевод в десятичную систему & Буква \\ \hline
0437 & $4 \cdot 16^2 + 3 \cdot 16^1 + 7 \cdot 16^0 = 1024 + 48 + 7 = 1079$ & з \\ \hline
0430 & $4 \cdot 16^2 + 3 \cdot 16^1 + 0 \cdot 16^0 = 1024 + 48 + 0 = 1072$ & а \\ \hline
043A & $4 \cdot 16^2 + 3 \cdot 16^1 + 10 \cdot 16^0 = 1024 + 48 + 10 = 1082$ & к \\ \hline
043E & $4 \cdot 16^2 + 3 \cdot 16^1 + 14 \cdot 16^0 = 1024 + 48 + 14 = 1086$ & о \\ \hline
043D & $4 \cdot 16^2 + 3 \cdot 16^1 + 13 \cdot 16^0 = 1024 + 48 + 13 = 1085$ & н \\ \hline
0430 & $4 \cdot 16^2 + 3 \cdot 16^1 + 0 \cdot 16^0 = 1024 + 48 + 0 = 1072$ & а \\ \hline
\end{tabular}
\end{center}

Все коды лежат в диапазоне 1072--1103, значит это строчные буквы; номер буквы равен код $- 1072$:
$1079 - 1072 = 7$ (з), $0$ (а), $10$ (к), $14$ (о), $13$ (н), $0$ (а).

Получено слово: \textbf{закона}.
\vspace{1cm}

\textbf{Отчет:}
\begin{enumerate}[nosep]
    \item Буквы русского алфавита и их коды в кодировках 866, 1251 и Unicode приведены в таблице выше.
    \item Фамилия: КОСЯКОВ.
    \item КОСЯКОВ (866): 138 142 145 159 138 142 130 (10) = 8A 8E 91 9F 8A 8E 82 (16).
    \item Отчество: Андреевич.
    \item Андреевич (1251): 192 237 228 240 229 229 226 232 247 (10) = C0 ED E4 F0 E5 E5 E2 E8 F7 (16).
    \item $\mathrm{04370430043A043E043D0430}_{16}$ = 1079 1072 1082 1086 1085 1072 (10) = <<закона>>.
\end{enumerate}

\end{document}
